\documentclass[a4paper,dvipdfmx]{jarticle}

\usepackage{graphicx}   % for graphics
\usepackage{mathtools}  % for math formulas
\usepackage{amssymb}    % for more math styles
\usepackage{amsfonts}    % for more math styles

\usepackage{xcolor}

\usepackage{hyperref}
\usepackage{pxjahyper}
\hypersetup{% hyperrefオプションリスト
 setpagesize=false,
 colorlinks=true,%
 linkcolor=blue,
 citecolor=red,
}





\renewcommand{\ttdefault}{pcr}

\usepackage{listings,jlisting}
\lstset{
 	language = Python,
 	basicstyle = \ttfamily\small,
 	commentstyle = \sffamily\gtfamily\color[gray]{.5},
 	keywordstyle = \bfseries,
 	framesep = 1em,
 	numbers = left,
 	stepnumber = 1,
 	numberstyle = \tiny\ttfamily,
    lineskip=-0.5ex,
    columns=[c]{fullflexible},
    breaklines=true,
    frame=tb,
}



% \pagestyle{empty}
% \thispagestyle{empty}



\usepackage[top=2cm, bottom=2cm, left=1cm, right=1cm, includefoot]{geometry}

\raggedbottom

\usepackage{bm}




\begin{document}

\title{機械学習}
\author{佐々木翔太}
\date{}

\maketitle

\section{分類}

\subsection{ロジスティック回帰による2値分類}

\subsubsection{目的}
ここでは，ロジスティック回帰を構成する各要素の必然性を明らかにする事を目的とする．

\subsubsection{2値分類タスクの問題設定}

2値分類タスク\footnote{以降は「2値分類」という言葉の使用を避け，「2値分類タスク」なのか「2値分類モデル」なのかを明示する．}とは，入力$\bm{x}$をその内容に応じて2種類に区別する問題である．クラスラベルは，2値の変数$k \in \{0, 1\}$で表現するものとする．2値分類タスクの目標は，入力$\bm{x}$から$k$を推定する事である．

この問題を定式化するため，入力$\bm{x}$を指定した時$k=1$となる事後確率\footnote{C. M. ビショップ：パターン認識と機械学習　上　ベイズ理論による統計的予測，p.16，シュプリンガー・ジャパン株式会社，東京（2008）．}$p(k=1|\bm{x})$を考える\footnote{$p$は，入力$\bm{x}$に対してダイナミックに変化するため，条件付き確率である．ここでは事後確率と呼ぶ．}．$\bm{x}$が与えられた時，事後確率$p(k=1|\bm{x})$を計算し，その値が$0.5$を超えれば$k=1$，下回れば$k=0$と判断する\footnote{岡谷貴之：深層学習　改訂第2版，p.19，講談社，東京（2024）．}．

ここでは，入力を$I$次元の列ベクトル$\bm{x} \in \mathbb{R}^I$と定義する．例えば画像の2値分類タスクであれば，$\bm{x}$は画像から抽出された単一の特徴ベクトルを意味する．また，セグメンテーションや物体検出などの2値分類タスクであっても，最終的には「特定のピクセルやグリッドにおける特徴ベクトルに対する独立した2値分類の反復」に帰着する\footnote{要出典}．よって，入力データの次元に依らず，2値分類モデルは「列ベクトル$\bm{x}$から事後確率$p \in (0, 1)$への写像（全射）\footnote{数学の景色：関数とは何か,写像とは何かを図解～定義と表記法と具体例～，https://mathlandscape.com/function/\#toc8，（2026年4月30日参照）．}」として一般化して議論できる．

\begin{align*}
  f: \mathbb{R}^I \to (0, 1) &,\quad \bm{x} \mapsto p \\
  \bm{x} = \begin{pmatrix}
    x_1 \\
    \vdots \\
    x_I
  \end{pmatrix} \in \mathbb{R}^I &,\quad p \in (0, 1)
\end{align*}

なお，$p \in (0, 1)$であり，$p \in [0, 1]$ではない．この理由は，ロジスティックシグモイド関数\footnote{「ロジスティック関数」とも呼ばれる．}の議論において後述する．

\subsubsection{2値分類モデルが満たすべき条件}

前項で定式化した2値分類モデルを設計するにあたり，モデルの出力値は少なくとも2つの条件を満たす必要がある．

\paragraph{条件1：不確実性を表現する「確率」である事}　%改行させるための全角スペース

現実の対象が持つ完全な情報を有限次元の特徴ベクトル$\bm{x}$に圧縮して観測する以上，観測されていない潜在変数\footnote{C. M. ビショップ：パターン認識と機械学習　下　ベイズ理論による統計的予測，p.76，シュプリンガー・ジャパン株式会社，東京（2008）．}\footnote{「隠れ変数」や「非観測変数」とも呼ばれる．}が存在する．この情報の欠落により，現実にはクラス$k=0$に属するデータと$k=1$に属するデータが，特徴空間上で全く同じ座標に位置する可能性がある．例えば，ある人物について「身長170cm」という特徴ベクトル$\bm{x}$が与えられ，その人物の性別を予測するタスクを考えると，正解ラベルは女性にも男性にもなり得る．

同一の入力$\bm{x}$に対して，正解ラベルが複数存在し得る以上，予測モデルは絶対的な$\{0, 1\}$の判断ではなく，不確実性を表す値を出力する必要がある．

そして，不確実性について合理的で一貫した推論を行う唯一の枠組みが確率論であるため，不確実性を表現するために確率を使う\footnote{C. M. ビショップ：パターン認識と機械学習　上　ベイズ理論による統計的予測，p.21，シュプリンガー・ジャパン株式会社，東京（2008）．}．

\paragraph{条件2：大部分において「微分可能な連続関数」である事}　%改行させるための全角スペース

前項により，モデルの出力値は確率となったが，これだけではモデルの学習は保証されない．ロジスティック回帰モデルのパラメータ最適化に勾配降下法を採用する事を仮定すると，入力$\bm{x}$から確率$p$への写像（モデル）は，定義域の大部分において微分可能な連続関数でなければならない\footnote{ReLUなどの活性化関数は特定の点において微分不可能である．しかし実用上は，劣微分を用いるなどして勾配降下法による学習が問題なく成立する．「大部分において」という表現は，このような実態を許容するものである．}．たとえ出力値が確率であっても，入力の変化に対して確率が階段状に変化する離散的な関数であった場合，その微分値は至る所で0となり，勾配が消失してしまう．よって学習を成立させるためには，モデルは単に確率を出力するだけでなく，入力に対して滑らかに連続変化する微分可能な関数として設計されなければならない．

\vskip\baselineskip
これまでの説明により，2値分類モデルの出力は「確率」であり，かつ「微分可能な連続関数」でなければならない事が示された．以降では，これらの条件を満たす具体的なモデルを設計する．

\subsubsection{モデル出力の設計と「ロジスティックシグモイド関数」の必然性}

機械学習の文脈で事後確率$p(k=1|\bm{x})$をモデル化するために，ニューラルネットワークを使用する．

\end{document}
